\chapter{Pemodelan dan Pelatihan}

Pada Bab 3, kita telah berhasil menyiapkan seluruh kebutuhan data untuk pemodelan. Mulai dari pembuatan subset 5 kelas daun tomat, penerapan transformasi gambar, hingga pembuatan \texttt{DataLoader} yang siap menyajikan data dalam bentuk \textit{batch}. Dataset dengan total 11.600 gambar (9.281 \textit{training} dan 2.319 \textit{validation}) kini telah siap untuk diproses oleh model.

Sekarang, tibalah fase yang paling inti dari praktikum ini: membangun dan melatih model \textit{Deep Learning}. Pada bab ini, kita akan membangun arsitektur model \textbf{ResNet9} menggunakan PyTorch, melatihnya dengan dataset daun tomat yang telah disiapkan, serta mengevaluasi seberapa baik model belajar mengenali pola-pola penyakit pada daun tomat.

Bab ini terbagi menjadi dua bagian utama: bagian \textbf{Pemodelan} (\textit{Modeling}) yang berfokus pada pembangunan arsitektur model, dan bagian \textbf{Pelatihan} (\textit{Training}) yang berfokus pada proses melatih model hingga menghasilkan performa yang optimal.

\section{Pemodelan (Modeling)}

Pada bagian ini, kita akan membangun arsitektur model lapis demi lapis. Model yang digunakan adalah \textbf{ResNet9}, yaitu varian arsitektur \textit{Residual Network} dengan 9 lapisan (\textit{layer}) yang dirancang khusus untuk klasifikasi gambar. ResNet9 menggunakan konsep \textit{residual connection} (koneksi sisa) yang memungkinkan model belajar lebih efektif dengan cara melewatkan informasi dari lapisan awal ke lapisan yang lebih dalam.

\subsection{Menyiapkan Device GPU/CPU}

Seperti yang telah disebutkan pada Bab 1, praktikum ini akan dijalankan di Google Colab yang menyediakan akses GPU. Penggunaan GPU dapat mempercepat proses \textit{training} model secara signifikan, terutama pada dataset gambar berukuran besar.

Pada tahap ini, kita membuat beberapa fungsi bantuan untuk memindahkan data dan model ke perangkat komputasi yang tersedia.

Jika GPU tersedia, proses training akan dijalankan menggunakan GPU. Jika tidak tersedia, proses training tetap dapat berjalan menggunakan CPU, tetapi waktu training akan lebih lama.

\begin{lstlisting}[language=Python, caption=Fungsi untuk menyiapkan device GPU/CPU]
# ============================================================
# FASE 4.1 - Menyiapkan Device GPU/CPU
# ============================================================

# ------------------------------------------------------------
# 1. Fungsi untuk memilih device
# ------------------------------------------------------------

def get_default_device():
    """
    Memilih GPU jika tersedia.
    Jika GPU tidak tersedia, maka menggunakan CPU.
    """
    if torch.cuda.is_available():
        return torch.device("cuda")
    else:
        return torch.device("cpu")


# ------------------------------------------------------------
# 2. Fungsi untuk memindahkan data ke device
# ------------------------------------------------------------

def to_device(data, device):
    """
    Memindahkan tensor atau sekumpulan tensor ke device yang dipilih.
    Device dapat berupa GPU atau CPU.
    """
    if isinstance(data, (list, tuple)):
        return [to_device(item, device) for item in data]

    return data.to(device, non_blocking=True)


# ------------------------------------------------------------
# 3. Class untuk memindahkan batch DataLoader ke device
# ------------------------------------------------------------

class DeviceDataLoader:
    """
    Membungkus DataLoader agar setiap batch data otomatis dipindahkan
    ke device yang digunakan, baik GPU maupun CPU.
    """

    def __init__(self, data_loader, device):
        self.data_loader = data_loader
        self.device = device

    def __iter__(self):
        """
        Menghasilkan batch data yang sudah dipindahkan ke device.
        """
        for batch in self.data_loader:
            yield to_device(batch, self.device)

    def __len__(self):
        """
        Mengembalikan jumlah batch dalam DataLoader.
        """
        return len(self.data_loader)


# ------------------------------------------------------------
# 4. Menentukan device yang digunakan
# ------------------------------------------------------------

device = get_default_device()

print("Device yang digunakan:", device)
\end{lstlisting}

Output yang akan tampil setelah kode dijalankan adalah sebagai berikut:

\begin{outputcode}[H]
\begin{lstlisting}[language={}]
Device yang digunakan: cuda
\end{lstlisting}
\caption{Output penentuan device komputasi}
\end{outputcode}

Fungsi \texttt{get\_default\_device()} akan mengecek ketersediaan GPU melalui \texttt{torch.cuda.is\_available()}. Jika GPU tersedia, fungsi mengembalikan \texttt{device='cuda'}. Jika tidak, fungsi mengembalikan \texttt{device='cpu'}.

Sementara itu, fungsi \texttt{to\_device()} digunakan untuk memindahkan data (\textit{tensor}) ke perangkat yang dipilih. Fungsi ini juga dapat menangani data dalam bentuk \textit{list} atau \textit{tuple} secara rekursif, sehingga kita tidak perlu memindahkan data satu per satu secara manual.

Selain kedua fungsi tersebut, kita juga membuat class \texttt{DeviceDataLoader} yang akan membungkus \texttt{DataLoader} dari Bab 3 agar setiap \textit{batch} data secara otomatis dipindahkan ke device yang digunakan. Dengan demikian, kita tidak perlu secara manual memindahkan data setiap kali akan digunakan untuk \textit{training} atau evaluasi.

\texttt{DeviceDataLoader} bekerja dengan cara membungkus \texttt{DataLoader} yang sudah ada. Ketika kita melakukan iterasi (\texttt{for batch in train\_dl}), setiap \textit{batch} data akan otomatis dipindahkan ke GPU atau CPU oleh fungsi \texttt{to\_device()} sebelum dikembalikan ke pemanggil.

\subsection{Memindahkan DataLoader ke Device}

Setelah device GPU/CPU berhasil ditentukan, langkah berikutnya adalah memindahkan data training dan validation ke device tersebut.

Jika GPU tersedia, setiap batch gambar dan label akan otomatis dipindahkan ke GPU agar proses training lebih cepat. Jika GPU tidak tersedia, data tetap diproses menggunakan CPU.

\begin{lstlisting}[language=Python, caption=Memindahkan DataLoader ke device]
# ============================================================
# FASE 4.2 - Memindahkan DataLoader ke Device
# ============================================================

device = get_default_device()

# Membungkus DataLoader hanya jika belum dibungkus
if not isinstance(train_dl, DeviceDataLoader):
    train_dl = DeviceDataLoader(train_dl, device)

if not isinstance(valid_dl, DeviceDataLoader):
    valid_dl = DeviceDataLoader(valid_dl, device)

print("Device yang digunakan:", device)
print("DataLoader siap digunakan pada device.")
\end{lstlisting}

Output yang akan tampil setelah kode dijalankan adalah sebagai berikut:

\begin{outputcode}[H]
\begin{lstlisting}[language={}]
Device yang digunakan: cuda
DataLoader siap digunakan pada device.
\end{lstlisting}
\caption{Output pemindahan DataLoader ke device}
\end{outputcode}

Sekarang, seluruh data yang akan digunakan dalam proses \textit{training} dan evaluasi sudah siap berada di perangkat komputasi yang sesuai. DataLoader yang telah dibungkus akan secara otomatis memindahkan setiap \textit{batch} gambar dan label ke device yang telah ditentukan.

Dengan \texttt{DeviceDataLoader}, kita tidak perlu lagi memanggil \texttt{.to(device)} secara manual setiap kali mengambil data dari DataLoader. Kode training menjadi lebih bersih dan efisien karena proses pemindahan device terjadi secara otomatis di balik layar.

\subsection{Membuat Residual Block Sederhana}

Pada tahap ini, kita membuat contoh sederhana dari Residual Block.

Residual Block adalah blok jaringan saraf yang menambahkan kembali input awal ke hasil proses konvolusi. Konsep ini membantu model mempertahankan informasi penting dari input dan membuat proses pembelajaran menjadi lebih stabil, terutama pada jaringan yang lebih dalam.

Pada contoh sederhana ini, input gambar diproses melalui dua layer konvolusi, kemudian hasilnya dijumlahkan kembali dengan input awal.

\begin{lstlisting}[language=Python, caption=Membuat SimpleResidualBlock]
# ============================================================
# FASE 4.3 - Membuat Residual Block Sederhana
# ============================================================

class SimpleResidualBlock(nn.Module):

    def __init__(self):
        super().__init__()

        self.conv1 = nn.Conv2d(
            in_channels=3,
            out_channels=3,
            kernel_size=3,
            stride=1,
            padding=1
        )

        self.relu1 = nn.ReLU()

        self.conv2 = nn.Conv2d(
            in_channels=3,
            out_channels=3,
            kernel_size=3,
            stride=1,
            padding=1
        )

        self.relu2 = nn.ReLU()

    def forward(self, x):
        # Proses konvolusi pertama
        out = self.conv1(x)
        out = self.relu1(out)

        # Proses konvolusi kedua
        out = self.conv2(out)

        # Menambahkan input awal ke output
        out = out + x

        # Aktivasi akhir
        out = self.relu2(out)

        return out
\end{lstlisting}

Penjelasan cara kerja \texttt{SimpleResidualBlock}:

\begin{itemize}
    \item Input gambar (\texttt{x}) diproses melalui lapisan konvolusi pertama (\texttt{conv1}) yang mempertahankan jumlah \textit{channel} dan dimensi spasial gambar.
    \item Hasil konvolusi pertama kemudian dilewatkan melalui fungsi aktivasi \texttt{ReLU} untuk menambahkan non-linearitas.
    \item Selanjutnya, data diproses melalui lapisan konvolusi kedua (\texttt{conv2}) dengan konfigurasi yang sama.
    \item Pada tahap inilah konsep \textit{residual connection} diterapkan --- hasil konvolusi kedua dijumlahkan dengan input awal (\texttt{out = out + x}). Penjumlahan ini memungkinkan informasi dari input awal tetap dipertahankan, sehingga model tidak kehilangan konteks penting meskipun telah melalui beberapa lapisan pemrosesan.
    \item Terakhir, hasil penjumlahan dilewatkan melalui aktivasi \texttt{ReLU} sekali lagi untuk menghasilkan output akhir.
\end{itemize}

Dengan adanya \textit{residual connection}, gradien dapat mengalir lebih mudah selama proses \textit{backpropagation}, sehingga model yang lebih dalam pun dapat dilatih secara efektif tanpa mengalami degradasi performa.

\subsection{Membuat Fungsi Akurasi dan Base Class Model}

Pada tahap ini, kita membuat fungsi untuk menghitung akurasi model dan class dasar untuk proses training serta validation.

Fungsi \texttt{accuracy()} digunakan untuk membandingkan hasil prediksi model dengan label asli. Sementara itu, \texttt{ImageClassificationBase} digunakan sebagai class dasar yang berisi alur training step, validation step, perhitungan rata-rata loss, rata-rata akurasi, dan tampilan hasil setiap epoch.

\begin{lstlisting}[language=Python, caption=Membuat fungsi akurasi dan base class]
# ============================================================
# FASE 5.4 - Membuat Fungsi Akurasi dan Base Class Model
# ============================================================

# ------------------------------------------------------------
# 1. Fungsi untuk menghitung akurasi
# ------------------------------------------------------------

def accuracy(outputs, labels):
    """
    Menghitung akurasi prediksi model.

    Parameter:
    outputs : hasil prediksi model sebelum softmax
    labels  : label asli dari dataset

    Return:
    nilai akurasi dalam bentuk tensor
    """

    # Mengambil index kelas dengan nilai prediksi tertinggi
    _, preds = torch.max(outputs, dim=1)

    # Membandingkan prediksi dengan label asli
    acc = torch.sum(preds == labels).item() / len(preds)

    return torch.tensor(acc)


# ------------------------------------------------------------
# 2. Base class untuk model klasifikasi gambar
# ------------------------------------------------------------

class ImageClassificationBase(nn.Module):
    """
    Class dasar untuk model klasifikasi gambar.

    Class ini berisi fungsi training, validation, dan logging hasil
    yang akan digunakan oleh model utama.
    """

    def training_step(self, batch):
        """
        Melakukan satu langkah training pada satu batch data.
        """

        images, labels = batch

        # Menghasilkan prediksi dari model
        outputs = self(images)

        # Menghitung loss training
        loss = F.cross_entropy(outputs, labels)

        return loss

    def validation_step(self, batch):
        """
        Melakukan satu langkah validation pada satu batch data.
        """

        images, labels = batch

        # Menghasilkan prediksi dari model
        outputs = self(images)

        # Menghitung loss validation
        loss = F.cross_entropy(outputs, labels)

        # Menghitung akurasi validation
        acc = accuracy(outputs, labels)

        return {
            "val_loss": loss.detach(),
            "val_accuracy": acc
        }

    def validation_epoch_end(self, outputs):
        """
        Menggabungkan hasil validation dari seluruh batch
        dalam satu epoch.
        """

        # Mengambil semua loss dari batch validation
        batch_losses = [x["val_loss"] for x in outputs]

        # Mengambil semua akurasi dari batch validation
        batch_accuracies = [x["val_accuracy"] for x in outputs]

        # Menghitung rata-rata loss dan akurasi
        epoch_loss = torch.stack(batch_losses).mean()
        epoch_accuracy = torch.stack(batch_accuracies).mean()

        return {
            "val_loss": epoch_loss,
            "val_accuracy": epoch_accuracy
        }

    def epoch_end(self, epoch, result):
        """
        Menampilkan hasil training dan validation pada akhir epoch.
        """

        print(
            "Epoch [{}], last_lr: {:.5f}, train_loss: {:.4f}, val_loss: {:.4f}, val_acc: {:.4f}".format(
                epoch,
                result["lrs"][-1],
                result["train_loss"],
                result["val_loss"],
                result["val_accuracy"]
            )
        )
\end{lstlisting}

Penjelasan masing-masing fungsi dan metode dalam \texttt{ImageClassificationBase}:

\begin{itemize}
    \item \textbf{\texttt{accuracy()}}: Fungsi ini mengambil indeks kelas dengan nilai tertinggi dari output model (\texttt{torch.max}), kemudian membandingkannya dengan label asli. Hasil perbandingan dibagi dengan jumlah total data untuk mendapatkan nilai akurasi dalam bentuk tensor.
    \item \textbf{\texttt{training\_step()}}: Menghitung \textit{loss} untuk satu \textit{batch} data \textit{training} menggunakan \texttt{cross\_entropy}, yang menggabungkan \textit{softmax} dan \textit{negative log likelihood} dalam satu langkah.
    \item \textbf{\texttt{validation\_step()}}: Menghitung \textit{loss} dan akurasi untuk satu \textit{batch} data \textit{validation}. Hasil \textit{loss} di-\textit{detach} agar tidak memengaruhi grafik komputasi.
    \item \textbf{\texttt{validation\_epoch\_end()}}: Menggabungkan seluruh hasil \textit{validation} dari setiap \textit{batch} menjadi satu nilai rata-rata untuk \textit{loss} dan akurasi.
    \item \textbf{\texttt{epoch\_end()}}: Menampilkan ringkasan hasil \textit{training} pada akhir setiap \textit{epoch}, termasuk \textit{learning rate} terakhir yang digunakan.
\end{itemize}

Perhatikan bahwa metode \texttt{epoch\_end} menggunakan key \texttt{result["lrs"]} yang berisi daftar \textit{learning rate} selama satu \textit{epoch}. Ini akan diisi oleh fungsi \textit{training} yang akan kita buat pada tahap selanjutnya.

\subsection{Membuat Arsitektur Model ResNet9}

Pada tahap ini, kita membuat arsitektur model yang akan digunakan untuk klasifikasi citra daun.

Model yang digunakan adalah \textbf{ResNet9}, yaitu arsitektur CNN sederhana yang menggunakan \textit{residual connection}. \textit{Residual connection} membantu model mempertahankan informasi penting dari layer sebelumnya, sehingga proses pembelajaran menjadi lebih stabil.

Model ini akan menerima gambar daun berukuran 256×256 piksel dengan 3 \textit{channel} warna RGB, kemudian menghasilkan output berupa 5 kelas kondisi daun tomat.

\begin{lstlisting}[language=Python, caption=Membangun arsitektur ResNet9]
# ============================================================
# FASE 5.5 - Membuat Arsitektur Model ResNet9
# ============================================================

# ------------------------------------------------------------
# 1. Membuat blok konvolusi
# ------------------------------------------------------------

def ConvBlock(in_channels, out_channels, pool=False):
    """
    Membuat blok konvolusi yang terdiri dari:
    - Conv2D
    - Batch Normalization
    - ReLU
    - MaxPooling opsional
    """

    layers = [
        nn.Conv2d(
            in_channels=in_channels,
            out_channels=out_channels,
            kernel_size=3,
            padding=1
        ),
        nn.BatchNorm2d(out_channels),
        nn.ReLU(inplace=True)
    ]

    if pool:
        layers.append(nn.MaxPool2d(4))

    return nn.Sequential(*layers)


# ------------------------------------------------------------
# 2. Membuat arsitektur ResNet9
# ------------------------------------------------------------

class ResNet9(ImageClassificationBase):
    """
    Arsitektur ResNet9 untuk klasifikasi citra daun.

    Input  : gambar RGB dengan 3 channel
    Output : jumlah kelas penyakit/kondisi daun
    """

    def __init__(self, in_channels, num_classes):
        super().__init__()

        # Blok awal
        self.conv1 = ConvBlock(in_channels, 64)

        # Downsampling pertama
        # Output: 128 x 64 x 64
        self.conv2 = ConvBlock(64, 128, pool=True)

        # Residual block pertama
        self.res1 = nn.Sequential(
            ConvBlock(128, 128),
            ConvBlock(128, 128)
        )

        # Downsampling kedua
        # Output: 256 x 16 x 16
        self.conv3 = ConvBlock(128, 256, pool=True)

        # Downsampling ketiga
        # Output: 512 x 4 x 4
        self.conv4 = ConvBlock(256, 512, pool=True)

        # Residual block kedua
        self.res2 = nn.Sequential(
            ConvBlock(512, 512),
            ConvBlock(512, 512)
        )

        # Classifier akhir
        self.classifier = nn.Sequential(
            nn.MaxPool2d(4),
            nn.Flatten(),
            nn.Linear(512, num_classes)
        )

    def forward(self, xb):
        """
        Forward pass model.
        """

        out = self.conv1(xb)
        out = self.conv2(out)

        # Residual connection pertama
        out = self.res1(out) + out

        out = self.conv3(out)
        out = self.conv4(out)

        # Residual connection kedua
        out = self.res2(out) + out

        out = self.classifier(out)

        return out
\end{lstlisting}

Penjelasan masing-masing komponen arsitektur ResNet9:

\begin{itemize}
    \item \textbf{\texttt{ConvBlock()}}: Fungsi pembantu yang membuat blok konvolusi standar, terdiri dari \texttt{Conv2D}, \texttt{BatchNorm2d}, dan \texttt{ReLU}. Jika parameter \texttt{pool=True}, sebuah lapisan \texttt{MaxPool2d(4)} ditambahkan untuk mereduksi dimensi spasial gambar.
    \item \textbf{\texttt{conv1}}: Blok konvolusi awal yang mengubah input 3 \textit{channel} (RGB) menjadi 64 \textit{channel} tanpa mengurangi dimensi.
    \item \textbf{\texttt{conv2}}: Blok konvolusi yang meningkatkan \textit{channel} dari 64 menjadi 128 dan mereduksi dimensi spasial menggunakan \texttt{MaxPool2d(4)}. Output berukuran 128×64×64.
    \item \textbf{\texttt{res1}}: \textit{Residual block} pertama yang terdiri dari dua blok konvolusi 128 \textit{channel}. Hasilnya dijumlahkan dengan input blok ini melalui operasi \texttt{self.res1(out) + out}. Penjumlahan inilah yang disebut sebagai \textit{residual connection}.
    \item \textbf{\texttt{conv3}}: Blok konvolusi yang meningkatkan \textit{channel} dari 128 menjadi 256 dengan reduksi dimensi. Output berukuran 256×16×16.
    \item \textbf{\texttt{conv4}}: Blok konvolusi yang meningkatkan \textit{channel} dari 256 menjadi 512. Output berukuran 512×4×4.
    \item \textbf{\texttt{res2}}: \textit{Residual block} kedua yang terdiri dari dua blok konvolusi 512 \textit{channel}, juga menggunakan \textit{residual connection}.
    \item \textbf{\texttt{classifier}}: Lapisan klasifikasi akhir yang terdiri dari \texttt{MaxPool2d(4)} untuk mereduksi dimensi menjadi 1×1, \texttt{Flatten} untuk mengubah matriks menjadi vektor, dan \texttt{Linear} yang memetakan 512 fitur menjadi 5 output kelas.
\end{itemize}

Arsitektur ini secara bertahap mengekstrak fitur-fitur visual dari gambar daun, mulai dari tepi dan tekstur sederhana di lapisan awal hingga pola penyakit yang kompleks di lapisan yang lebih dalam. Penggunaan dua \textit{residual connection} memastikan bahwa informasi penting dari lapisan sebelumnya tetap dipertahankan selama proses pemrosesan.

\subsection{Membuat Object Model ResNet9}

Setelah arsitektur ResNet9 berhasil dibuat, langkah berikutnya adalah membuat object model.

Model dibuat dengan input 3 channel karena gambar yang digunakan adalah gambar RGB. Jumlah output model disesuaikan dengan jumlah kelas pada dataset training. Pada prototype ini, jumlah kelas yang digunakan adalah 5 kelas kondisi daun tomat.

Setelah model dibuat, model dipindahkan ke device yang digunakan, yaitu GPU jika tersedia atau CPU jika GPU tidak tersedia.

\begin{lstlisting}[language=Python, caption=Membuat object model ResNet9]
# ============================================================
# FASE 5.6 - Membuat Object Model ResNet9
# ============================================================

# Jumlah channel input gambar RGB
in_channels = 3

# Jumlah kelas output sesuai dataset training
num_classes = len(train.classes)

# Membuat model ResNet9
model = ResNet9(in_channels, num_classes)

# Memindahkan model ke device GPU/CPU
model = to_device(model, device)

# Menampilkan model
print("Model ResNet9 berhasil dibuat.")
print("Device yang digunakan :", device)
print("Jumlah kelas output   :", num_classes)
print("Daftar kelas          :", train.classes)

model
\end{lstlisting}

Output yang akan tampil setelah kode dijalankan adalah sebagai berikut:

\begin{outputcode}[H]
\begin{lstlisting}[language={}]
Model ResNet9 berhasil dibuat.
Device yang digunakan : cuda
Jumlah kelas output   : 5
Daftar kelas          : ['Bacterial Spot', 'Early Blight', 'Healthy', 'Late Blight', 'Leaf Mold']
ResNet9(
  (conv1): Sequential(
    (0): Conv2d(3, 64, kernel_size=(3, 3), stride=(1, 1), padding=(1, 1))
    (1): BatchNorm2d(64, eps=1e-05, momentum=0.1, affine=True, track_running_stats=True)
    (2): ReLU(inplace=True)
  )
  (conv2): Sequential(
    (0): Conv2d(64, 128, kernel_size=(3, 3), stride=(1, 1), padding=(1, 1))
    (1): BatchNorm2d(128, eps=1e-05, momentum=0.1, affine=True, track_running_stats=True)
    (2): ReLU(inplace=True)
    (3): MaxPool2d(kernel_size=4, stride=4, padding=0, dilation=1, ceil_mode=False)
  )
  (res1): Sequential(
    (0): Sequential(
      (0): Conv2d(128, 128, kernel_size=(3, 3), stride=(1, 1), padding=(1, 1))
      (1): BatchNorm2d(128, eps=1e-05, momentum=0.1, affine=True, track_running_stats=True)
      (2): ReLU(inplace=True)
    )
    (1): Sequential(
      (0): Conv2d(128, 128, kernel_size=(3, 3), stride=(1, 1), padding=(1, 1))
      (1): BatchNorm2d(128, eps=1e-05, momentum=0.1, affine=True, track_running_stats=True)
      (2): ReLU(inplace=True)
    )
  )
  (conv3): Sequential(
    (0): Conv2d(128, 256, kernel_size=(3, 3), stride=(1, 1), padding=(1, 1))
    (1): BatchNorm2d(256, eps=1e-05, momentum=0.1, affine=True, track_running_stats=True)
    (2): ReLU(inplace=True)
    (3): MaxPool2d(kernel_size=4, stride=4, padding=0, dilation=1, ceil_mode=False)
  )
  (conv4): Sequential(
    (0): Conv2d(256, 512, kernel_size=(3, 3), stride=(1, 1), padding=(1, 1))
    (1): BatchNorm2d(512, eps=1e-05, momentum=0.1, affine=True, track_running_stats=True)
    (2): ReLU(inplace=True)
    (3): MaxPool2d(kernel_size=4, stride=4, padding=0, dilation=1, ceil_mode=False)
  )
  (res2): Sequential(
    (0): Sequential(
      (0): Conv2d(512, 512, kernel_size=(3, 3), stride=(1, 1), padding=(1, 1))
      (1): BatchNorm2d(512, eps=1e-05, momentum=0.1, affine=True, track_running_stats=True)
      (2): ReLU(inplace=True)
    )
    (1): Sequential(
      (0): Conv2d(512, 512, kernel_size=(3, 3), stride=(1, 1), padding=(1, 1))
      (1): BatchNorm2d(512, eps=1e-05, momentum=0.1, affine=True, track_running_stats=True)
      (2): ReLU(inplace=True)
    )
  )
  (classifier): Sequential(
    (0): MaxPool2d(kernel_size=4, stride=4, padding=0, dilation=1, ceil_mode=False)
    (1): Flatten(start_dim=1, end_dim=-1)
    (2): Linear(in_features=512, out_features=5, bias=True)
  )
)
\end{lstlisting}
\caption{Output struktur model ResNet9}
\end{outputcode}

Dari output di atas, kita dapat melihat bahwa model ResNet9 telah berhasil dibuat dengan struktur yang sesuai dengan arsitektur yang telah didefinisikan sebelumnya. Lapisan terakhir (\texttt{Linear(512, 5)}) mengonfirmasi bahwa model akan menghasilkan 5 output, sesuai dengan jumlah kelas daun tomat yang telah ditetapkan sejak Bab 1.

Informasi tambahan yang ditampilkan mencakup device yang digunakan (GPU/CPU), jumlah kelas output, dan daftar nama kelas yang akan diprediksi oleh model, yaitu \textbf{Bacterial Spot}, \textbf{Early Blight}, \textbf{Healthy}, \textbf{Late Blight}, dan \textbf{Leaf Mold}.

\subsection{Menampilkan Ringkasan Arsitektur Model}

Setelah model berhasil dibuat, langkah berikutnya adalah menampilkan ringkasan arsitektur model.

Ringkasan ini digunakan untuk melihat susunan layer, ukuran output setiap layer, jumlah parameter, serta estimasi kebutuhan memori model. Tahap ini penting untuk memastikan bahwa model sudah sesuai dengan input gambar berukuran 256×256 piksel dan menghasilkan output sesuai jumlah kelas yang digunakan.

\begin{lstlisting}[language=Python, caption=Menampilkan ringkasan arsitektur model]
# ============================================================
# FASE 5.7 - Menampilkan Ringkasan Arsitektur Model
# ============================================================

from torchsummary import summary

# Menentukan ukuran input gambar
INPUT_SHAPE = (3, 256, 256)

# Mengecek device model
model_device = next(model.parameters()).device

# torchsummary hanya menerima "cuda" atau "cpu"
summary_device = "cuda" if model_device.type == "cuda" else "cpu"

print("Input shape model :", INPUT_SHAPE)
print("Device model      :", model_device)
print("Device summary    :", summary_device)
print("Jumlah kelas      :", len(train.classes))

# Menampilkan ringkasan arsitektur model
summary(model, input_size=INPUT_SHAPE, device=summary_device)
\end{lstlisting}

Output yang akan tampil setelah kode dijalankan adalah sebagai berikut:

\begin{outputcode}[H]
\begin{lstlisting}[language={}]
Input shape model : (3, 256, 256)
Device model      : cuda:0
Device summary    : cuda
Jumlah kelas      : 5
----------------------------------------------------------------
        Layer (type)               Output Shape         Param #
================================================================
            Conv2d-1         [-1, 64, 256, 256]           1,792
       BatchNorm2d-2         [-1, 64, 256, 256]             128
              ReLU-3         [-1, 64, 256, 256]               0
            Conv2d-4        [-1, 128, 256, 256]          73,856
       BatchNorm2d-5        [-1, 128, 256, 256]             256
              ReLU-6        [-1, 128, 256, 256]               0
         MaxPool2d-7          [-1, 128, 64, 64]               0
            Conv2d-8          [-1, 128, 64, 64]         147,584
       BatchNorm2d-9          [-1, 128, 64, 64]             256
             ReLU-10          [-1, 128, 64, 64]               0
           Conv2d-11          [-1, 128, 64, 64]         147,584
      BatchNorm2d-12          [-1, 128, 64, 64]             256
             ReLU-13          [-1, 128, 64, 64]               0
           Conv2d-14          [-1, 256, 64, 64]         295,168
      BatchNorm2d-15          [-1, 256, 64, 64]             512
             ReLU-16          [-1, 256, 64, 64]               0
        MaxPool2d-17          [-1, 256, 16, 16]               0
           Conv2d-18          [-1, 512, 16, 16]       1,180,160
      BatchNorm2d-19          [-1, 512, 16, 16]           1,024
             ReLU-20          [-1, 512, 16, 16]               0
        MaxPool2d-21            [-1, 512, 4, 4]               0
           Conv2d-22            [-1, 512, 4, 4]       2,359,808
      BatchNorm2d-23            [-1, 512, 4, 4]           1,024
             ReLU-24            [-1, 512, 4, 4]               0
           Conv2d-25            [-1, 512, 4, 4]       2,359,808
      BatchNorm2d-26            [-1, 512, 4, 4]           1,024
             ReLU-27            [-1, 512, 4, 4]               0
        MaxPool2d-28            [-1, 512, 1, 1]               0
          Flatten-29                  [-1, 512]               0
           Linear-30                    [-1, 5]           2,565
================================================================
Total params: 6,572,805
Trainable params: 6,572,805
Non-trainable params: 0
----------------------------------------------------------------
Input size (MB): 0.75
Forward/backward pass size (MB): 343.95
Params size (MB): 25.07
Estimated Total Size (MB): 369.77
----------------------------------------------------------------
\end{lstlisting}
\caption{Output ringkasan arsitektur model ResNet9}
\end{outputcode}

Dari ringkasan di atas, terdapat beberapa informasi penting:

\begin{itemize}
    \item \textbf{Total parameter}: Model memiliki \textbf{6.572.805} parameter yang semuanya akan dipelajari selama proses \textit{training}. Jumlah ini tergolong rendah untuk arsitektur \textit{deep learning}, sehingga model dapat dilatih dengan relatif cepat.
    \item \textbf{Input size}: Setiap gambar input berukuran 0.75 MB (3 \textit{channel} × 256 × 256 piksel).
    \item \textbf{Estimated Total Size}: Diperkirakan total memori yang dibutuhkan adalah sekitar \textbf{369.77 MB} untuk satu kali \textit{forward} dan \textit{backward pass}. Jumlah ini masih dalam batas wajar untuk GPU yang tersedia di Google Colab (biasanya 12-16 GB).
    \item Lapisan \texttt{Linear-30} menunjukkan bahwa output akhir adalah vektor dengan 5 nilai, yang sesuai dengan jumlah kelas daun tomat yang digunakan dalam praktikum ini.
\end{itemize}

Dengan demikian, arsitektur model ResNet9 telah selesai dibangun dan siap untuk memasuki fase pelatihan.

\section{Pelatihan (Training)}

Setelah model berhasil dibangun, langkah selanjutnya adalah melatih model menggunakan data \textit{training} yang telah disiapkan. Pada bagian ini, kita akan membuat fungsi-fungsi yang diperlukan untuk proses \textit{training}, mengevaluasi performa awal model, menentukan \textit{hyperparameter}, dan menjalankan proses \textit{training} secara penuh.

\subsection{Membuat Fungsi Training dan Evaluasi Model}

Pada tahap ini, kita membuat beberapa fungsi utama untuk proses training.

Fungsi \texttt{evaluate()} digunakan untuk mengevaluasi model pada data validation. Fungsi \texttt{get\_lr()} digunakan untuk mengambil nilai learning rate saat training berlangsung. Fungsi \texttt{fit\_one\_cycle()} digunakan untuk menjalankan proses training model menggunakan strategi \textbf{One Cycle Learning Rate}.

Strategi \textit{One Cycle Learning Rate} adalah strategi pengaturan learning rate yang naik dan turun selama proses training agar model dapat belajar lebih stabil. Learning rate dimulai dari nilai yang rendah, naik secara bertahap hingga mencapai nilai maksimum, kemudian turun kembali menjelang akhir training.

\begin{lstlisting}[language=Python, caption=Membuat fungsi training dan evaluasi model]
# ============================================================
# FASE 5.1 - Membuat Fungsi Training dan Evaluasi Model
# ============================================================

# ------------------------------------------------------------
# 1. Fungsi evaluasi model pada data validation
# ------------------------------------------------------------

@torch.no_grad()
def evaluate(model, val_loader):
    """
    Mengevaluasi model menggunakan data validation.

    Fungsi ini tidak menghitung gradient karena hanya digunakan
    untuk evaluasi, bukan untuk training.
    """

    # Mengubah model ke mode evaluasi
    model.eval()

    # Menyimpan hasil validation dari setiap batch
    outputs = [
        model.validation_step(batch)
        for batch in val_loader
    ]

    # Menggabungkan hasil seluruh batch validation
    return model.validation_epoch_end(outputs)


# ------------------------------------------------------------
# 2. Fungsi untuk mengambil nilai learning rate
# ------------------------------------------------------------

def get_lr(optimizer):
    """
    Mengambil nilai learning rate dari optimizer.
    """

    for param_group in optimizer.param_groups:
        return param_group["lr"]


# ------------------------------------------------------------
# 3. Fungsi training dengan One Cycle Learning Rate
# ------------------------------------------------------------

def fit_one_cycle(
    epochs,
    max_lr,
    model,
    train_loader,
    val_loader,
    weight_decay=0,
    grad_clip=None,
    opt_func=torch.optim.Adam
):
    """
    Melatih model menggunakan One Cycle Learning Rate.

    Parameter:
    epochs       : jumlah epoch training
    max_lr       : learning rate maksimum
    model        : model yang akan dilatih
    train_loader : DataLoader untuk data training
    val_loader   : DataLoader untuk data validation
    weight_decay : regularisasi untuk mengurangi overfitting
    grad_clip    : batas maksimum gradient
    opt_func     : optimizer yang digunakan
    """

    # Membersihkan cache GPU jika tersedia
    if torch.cuda.is_available():
        torch.cuda.empty_cache()

    # Menyimpan history training
    history = []

    # Membuat optimizer
    optimizer = opt_func(
        model.parameters(),
        lr=max_lr,
        weight_decay=weight_decay
    )

    # Scheduler One Cycle Learning Rate
    scheduler = torch.optim.lr_scheduler.OneCycleLR(
        optimizer,
        max_lr=max_lr,
        epochs=epochs,
        steps_per_epoch=len(train_loader)
    )

    # Proses training per epoch
    for epoch in range(epochs):

        # Mengubah model ke mode training
        model.train()

        # Menyimpan loss training dan learning rate
        train_losses = []
        lrs = []

        # Proses training per batch
        for batch in train_loader:

            # Menghitung loss training
            loss = model.training_step(batch)

            # Menyimpan loss
            train_losses.append(loss.detach())

            # Backpropagation
            loss.backward()

            # Gradient clipping jika digunakan
            if grad_clip is not None:
                nn.utils.clip_grad_value_(
                    model.parameters(),
                    grad_clip
                )

            # Update parameter model
            optimizer.step()

            # Reset gradient
            optimizer.zero_grad()

            # Menyimpan learning rate saat ini
            lrs.append(get_lr(optimizer))

            # Update learning rate
            scheduler.step()

        # Evaluasi model pada data validation
        result = evaluate(model, val_loader)

        # Menambahkan rata-rata training loss
        result["train_loss"] = torch.stack(train_losses).mean().item()

        # Menambahkan daftar learning rate
        result["lrs"] = lrs

        # Menampilkan hasil epoch
        model.epoch_end(epoch, result)

        # Menyimpan hasil epoch ke history
        history.append(result)

    return history
\end{lstlisting}

Penjelasan masing-masing fungsi:

\begin{itemize}
    \item \textbf{\texttt{evaluate()}}: Fungsi ini menggunakan dekorator \texttt{@torch.no\_grad()} agar proses evaluasi tidak menyimpan grafik komputasi, sehingga lebih hemat memori. Model di-set ke mode \texttt{eval()} untuk menonaktifkan \textit{Batch Normalization} yang bersifat stokastik. Fungsi ini memanggil \texttt{model.validation\_step()} untuk setiap \textit{batch}, kemudian menggabungkan hasilnya melalui \texttt{model.validation\_epoch\_end()}.
    \item \textbf{\texttt{get\_lr()}}: Fungsi sederhana untuk mengambil nilai \textit{learning rate} yang sedang digunakan oleh \textit{optimizer}. Nilai ini akan dicatat pada setiap \textit{batch} untuk dipantau perkembangannya.
    \item \textbf{\texttt{fit\_one\_cycle()}}: Fungsi utama yang mengoordinasikan seluruh proses \textit{training}. Fungsi ini menginisialisasi \textit{optimizer} (Adam) dan \textit{One Cycle Learning Rate Scheduler}, kemudian menjalankan loop \textit{training} untuk setiap \textit{epoch}. Pada setiap \textit{epoch}, model melakukan:\begin{itemize}
        \item \textit{Training} per \textit{batch} menggunakan \texttt{model.training\_step()}.
        \item \textit{Backpropagation} untuk menghitung gradien.
        \item \textit{Gradient clipping} (jika dikonfigurasi) untuk mencegah gradien meledak.
        \item Pembaruan bobot model melalui \textit{optimizer}.
        \item Pencatatan \textit{learning rate} dan pembaruan \textit{scheduler}.
        \item Evaluasi performa model pada data \textit{validation}.
    \end{itemize}
\end{itemize}

\subsection{Evaluasi Awal Model Sebelum Training}

Sebelum proses training dimulai, model dievaluasi terlebih dahulu menggunakan data validation.

Tahap ini bertujuan untuk melihat performa awal model sebelum belajar dari data training. Karena model belum dilatih, nilai akurasi biasanya masih rendah dan loss masih tinggi.

Hasil evaluasi awal ini akan disimpan ke dalam variabel \texttt{history} agar nantinya dapat dibandingkan dengan hasil setelah training.

\begin{lstlisting}[language=Python, caption=Evaluasi awal model sebelum training]
# ============================================================
# FASE 5.2 - Evaluasi Awal Model Sebelum Training
# ============================================================

%%time

# Mengevaluasi model sebelum training
initial_result = evaluate(model, valid_dl)

# Menyimpan hasil evaluasi awal ke dalam history
history = [initial_result]

# Menampilkan hasil evaluasi awal
history
\end{lstlisting}

Output yang akan tampil setelah kode dijalankan adalah sebagai berikut:

\begin{outputcode}[H]
\begin{lstlisting}[language={}]
CPU times: user 7.61 s, sys: 559 ms, total: 8.17 s
Wall time: 9.58 s
[{'val_loss': tensor(1.6104, device='cuda:0'), 'val_accuracy': tensor(0.2277)}]
\end{lstlisting}
\caption{Output evaluasi awal model}
\end{outputcode}

Hasil evaluasi awal menunjukkan bahwa model memiliki \textit{validation loss} sebesar \textbf{1.6104} dan akurasi sebesar \textbf{22.77\%} --- nilai yang sedikit di atas tebakan acak (20\% untuk 5 kelas). Hal ini wajar karena model masih memiliki bobot acak dan belum pernah melihat data training sebelumnya.

Perhatikan penggunaan \texttt{\%\%time} pada sel kode di atas. Ini adalah \textit{magic command} khusus Jupyter Notebook (Google Colab) yang akan menampilkan waktu eksekusi sel. Informasi waktu ini berguna untuk memperkirakan berapa lama proses training nantinya akan berlangsung.

Hasil evaluasi disimpan ke dalam \texttt{history = [initial\_result]} sebagai elemen pertama dari sebuah \textit{list}. Nantinya, setelah proses training selesai, \texttt{history} akan berisi hasil dari setiap \textit{epoch} dan dapat digunakan untuk membandingkan perkembangan performa model dari waktu ke waktu.

\subsection{Menentukan Hyperparameter Training}

Sebelum proses training dijalankan, kita perlu menentukan beberapa hyperparameter utama.

Hyperparameter adalah nilai konfigurasi yang mengatur bagaimana model belajar, seperti jumlah epoch, learning rate, gradient clipping, weight decay, dan optimizer.

Pada prototype ini, jumlah epoch dibuat kecil terlebih dahulu agar proses training lebih cepat saat uji coba. Jika hasil model sudah stabil, jumlah epoch dapat ditingkatkan.

Berikut adalah \textit{hyperparameter} yang akan digunakan dalam praktikum ini:

\begin{itemize}
    \item \textbf{Epochs}: Jumlah iterasi penuh model melihat seluruh data \textit{training}. Semakin banyak \textit{epoch}, semakin lama model belajar. Pada praktikum ini, kita menggunakan \textbf{2 epoch}.
    \item \textbf{Max LR} (\textit{Maximum Learning Rate}): Nilai \textit{learning rate} tertinggi yang akan digunakan oleh \textit{One Cycle Scheduler}. Nilai \textbf{0.01} dipilih agar model dapat belajar dengan kecepatan yang optimal.
    \item \textbf{Grad Clip}: Batas maksimum nilai gradien untuk mencegah gradien meledak (\textit{exploding gradient}). Nilai \textbf{0.1} digunakan untuk menjaga stabilitas \textit{training}.
    \item \textbf{Weight Decay}: Teknik regularisasi yang menambahkan penalti pada bobot model agar tidak terlalu besar. Nilai \textbf{1e-4} digunakan untuk mencegah \textit{overfitting}.
    \item \textbf{Optimizer}: Algoritma yang digunakan untuk memperbarui bobot model. \textbf{Adam} (\textit{Adaptive Moment Estimation}) dipilih karena kemampuannya mengatur \textit{learning rate} secara adaptif untuk setiap parameter.
\end{itemize}

\begin{lstlisting}[language=Python, caption=Menentukan hyperparameter training]
# ============================================================
# FASE 5.3 - Menentukan Hyperparameter Training
# ============================================================

# Jumlah epoch training
epochs = 2

# Learning rate maksimum untuk One Cycle Learning Rate
max_lr = 0.01

# Batas maksimum gradient untuk mencegah gradient terlalu besar
grad_clip = 0.1

# Regularisasi untuk membantu mengurangi overfitting
weight_decay = 1e-4

# Optimizer yang digunakan
opt_func = torch.optim.Adam

# Menampilkan konfigurasi training
print("Konfigurasi Training")
print("Jumlah epoch      :", epochs)
print("Max learning rate :", max_lr)
print("Gradient clipping :", grad_clip)
print("Weight decay      :", weight_decay)
print("Optimizer         :", opt_func.__name__)
\end{lstlisting}

Output yang akan tampil setelah kode dijalankan adalah sebagai berikut:

\begin{outputcode}[H]
\begin{lstlisting}[language={}]
Konfigurasi Training
Jumlah epoch      : 2
Max learning rate : 0.01
Gradient clipping : 0.1
Weight decay      : 0.0001
Optimizer         : Adam
\end{lstlisting}
\caption{Output konfigurasi hyperparameter}
\end{outputcode}

\subsection{Menjalankan Training Model}

Setelah hyperparameter ditentukan, langkah berikutnya adalah menjalankan proses training model.

Pada tahap ini, model dilatih menggunakan data training dan dievaluasi menggunakan data validation pada setiap epoch. Hasil training seperti \texttt{train\_loss}, \texttt{val\_loss}, \texttt{val\_accuracy}, dan learning rate akan disimpan ke dalam variabel \texttt{history}.

Variabel \texttt{history} nantinya digunakan untuk membuat visualisasi perkembangan akurasi, loss, dan learning rate model.

\begin{lstlisting}[language=Python, caption=Menjalankan training model]
%%time

# ============================================================
# FASE 5.4 - Menjalankan Training Model
# ============================================================

# Menjalankan training model menggunakan One Cycle Learning Rate
training_history = fit_one_cycle(
    epochs=epochs,
    max_lr=max_lr,
    model=model,
    train_loader=train_dl,
    val_loader=valid_dl,
    grad_clip=grad_clip,
    weight_decay=weight_decay,
    opt_func=opt_func
)

# Menambahkan hasil training ke history
history += training_history

# Menampilkan history training
history
\end{lstlisting}

Output yang akan tampil setelah kode dijalankan adalah sebagai berikut:

\begin{outputcode}[H]
\begin{lstlisting}[language={}]
Epoch [0], last_lr: 0.00812, train_loss: 0.7888, val_loss: 5.3282, val_acc: 0.4178
Epoch [1], last_lr: 0.00000, train_loss: 0.2676, val_loss: 0.1044, val_acc: 0.9628
CPU times: user 2min 15s, sys: 1min 53s, total: 4min 8s
Wall time: 4min 18s
\end{lstlisting}
\caption{Output proses training model}
\end{outputcode}

Setelah training selesai, variabel \texttt{history} kini berisi tiga elemen: hasil evaluasi awal (dari tahap sebelumnya) ditambah hasil dua epoch training. Variabel ini dapat ditampilkan untuk melihat detail lengkap, termasuk nilai \textit{learning rate} untuk setiap \textit{batch}:

\begin{outputcode}[H]
\begin{lstlisting}[language={}]
[{'val_loss': tensor(1.6104, device='cuda:0'), 'val_accuracy': tensor(0.2277)},
 {'val_loss': tensor(5.3282, device='cuda:0'),
  'val_accuracy': tensor(0.4178),
  'train_loss': 0.7888427972793579,
  'lrs': [0.0004, 0.0004, ..., 0.0081, 0.0081]},
 {'val_loss': tensor(0.1044, device='cuda:0'),
  'val_accuracy': tensor(0.9628),
  'train_loss': 0.2675580680370331,
  'lrs': [0.0081, 0.0081, ..., 6e-07, 4e-08]}]
\end{lstlisting}
\caption{Output history training (LR list diringkas)}
\end{outputcode}

Berdasarkan hasil training, model menunjukkan perkembangan performa yang sangat baik.

Sebelum training, model memiliki \texttt{val\_loss} sebesar \textbf{1.6104} dan \texttt{val\_accuracy} sebesar \textbf{0.2277} atau sekitar \textbf{22.77\%}. Nilai ini wajar karena model belum belajar dari data training, sehingga prediksinya masih mendekati tebakan acak. Karena dataset memiliki 5 kelas, tebakan acak memang berada di sekitar 20\%.

Setelah training epoch pertama, nilai \texttt{train\_loss} menjadi \textbf{0.7888}, \texttt{val\_loss} menjadi \textbf{5.3282}, dan \texttt{val\_accuracy} naik menjadi \textbf{0.4178} atau sekitar \textbf{41.78\%}. Hal ini menunjukkan bahwa model mulai berhasil mempelajari pola visual dari citra daun, meskipun \texttt{val\_loss} masih cukup tinggi.

Pada epoch kedua, performa model meningkat lebih jauh. Nilai \texttt{train\_loss} turun menjadi \textbf{0.2676}, \texttt{val\_loss} turun drastis menjadi \textbf{0.1044}, dan \texttt{val\_accuracy} naik menjadi \textbf{0.9628} atau sekitar \textbf{96.28\%}. Artinya, model mampu mengklasifikasikan sebagian besar data validation dengan benar.

Secara keseluruhan, hasil ini menunjukkan bahwa model ResNet9 berhasil belajar dengan baik dari dataset lima kelas daun tomat, yaitu \textbf{Bacterial Spot}, \textbf{Early Blight}, \textbf{Healthy}, \textbf{Late Blight}, dan \textbf{Leaf Mold}.
